\section{Nine-variable equations: positive signed kernels and tunable norms}
The new source encoding and the free-carrier argument combine without any
additional unknowns.  For positive integer unknowns
\(b,g,h,s,w,C,i,j,n\), the following equations have no positive integral
zero if and only if \(\operatorname{Con}(\mathrm{ZFC})\), provably in ZFC,
using the inherited source/compiler and numerical ratio theorems.
\begin{center}
\begin{tabular}{r|r|r}
Majorant order & Exact degree & Binary arithmetic operations\\\hline
Baseline & 1,731,700 & 462\\
1 & 1,663,788 & 475\\
2 & 1,595,876 & 478\\
3 & 1,527,964 & 482\\
4 & 1,494,006 & 487
\end{tabular}
\end{center}
The baseline improves both measures over the previous
\(1{,}910{,}404/473\) construction.  The other rows are degree/cost
tradeoffs, not claims to dominate on every measure.  An order-zero variant
has degree \(1{,}731{,}700\) and 465 operations and is retained only for checking.

\subsection{The source and numerical interface}
Use the unchanged original31 four-coefficient tuple, supplied as binary
payloads, and the signed split kernels of the preceding construction.
Their degrees are \(1939,1842,1738\), with positive leading coefficients
\(2,2,\beta\).  For \(M=2^{2134927}\), every coefficient has absolute
value less than \(M\), so each kernel is positive when \(B>2M\).
Let \(\kappa=2134940\), \(a=2^\kappa\), \(B=ab^2\), and define
\begin{align*}
 \ell&=\sum_i B^{w_i},& U_s&=\ell+g,&
 G_s&=(D_2(B)U_s+D_1(B))U_s+D_0(B),\\
 Q_s&=\sum_j B^{q_j},& N_0&=B^{109},&
 T_s&=N_0-1-(b-1)\ell+N_0(B-2)Q_s,\\
 S_s&=g+N_0(2G_s+Q_s),& L&=bB^{2064},&
 N&=2^{16012039}B^{1090},\\
 R&=L(S_s+T_s+1)-S_s-1,&t&=bB^{108}.&&
\end{align*}
The maximum weight is 108, convolution height is 1954, and the maximum
target is 1939.  The global sign reversal preserves the source zero set.
The finite source certificate verifies all 22,544 isolation interactions.
For binary \(b=2^k,B=2^v\), the calibrated threshold is
\[
 \operatorname{pop}(R)=2195v-30k-22-\delta
                      =2\log_2 N-\delta,
\]
where \(\delta\ge0\) is precisely the carry defect.  The even value of
\(\kappa\) matters: \(N\) has no separate factor \(b\).

Set \(j_0=j-1\) and define the shared ratio registers
\begin{align*}
 U&=N^2w,&Y&=N^2s,&M_0&=RY,\\
 P&=2M_0^2U,&K&=R+1+h(P-1),& A&=M_0(U+1),\\
 B_1&=2R+1,&W_1&=bw,&J&=4A-5,\\
 Q&=3W_1C-2(W_1^2-1),&D&=(A^2-1)C^2+1,&e&=iJC,\\
 E&=eDC,&F&=(A^2-1)E^2+1,&G&=A+F(CD+1-A),\\
 H&=C+F(B_1+j_0C),&I&=(G^2-1)H^2+1,&&\\
 a_0&=(P^2-1)K^2+1,&a_1&=DFI,&V_0&=PK,\\
 m&=K^2(4t-g)-16t(C-KY)^2,&\Omega&=Jm,&E_0&=J(n+W_1^2)+Q.
\end{align*}
The original numerical interface becomes exactly: \(a_0,a_1\) are
squares, \(J\mid Q\), and \(m>0\).  The latter recovers
\(C>K(Y-1/2)>K\) and \(C>B_1+W_1\); it is not an assumed domain restriction.
The exact old-index correspondence is
\(j_{\rm old}=1+((F-1)/C)B_1+Fj_0\).

\subsection{Higher majorants without extra congruence restrictions}
The baseline takes \(\lambda=8\) and
\(T=e(8D^2-4D-1)G\).  For order \(r=0,1,2,3,4\), put
\(s_r=128,64,8,4,1\), respectively, and \(\lambda=2048/s_r\).
Define
\begin{align*}
 L_m&=CD+1-A,&N_*&=2G+L_m,\\
 N_1&=2e^2D(D-1)L_m+3C,&N_2&=2e^2(D-1)L_m,\\
 N_3&=2e^2(L_m-C),&N_4&=2e^2C.
\end{align*}
With \((c_1,c_2,c_3,c_4)=(64,40,28,21)\), let
\[
 T=e\left[\frac{1024D^2-512D-128}{s_r}N_*
       -\sum_{j=1}^r\frac{c_j}{s_r}N_j\right]+\epsilon_r,
 \qquad V_1=TH,
\]
where \(\epsilon_0=\lambda\),
\(\epsilon_r=512i/s_r+\lambda\) for \(1\le r\le3\), and
\(\epsilon_4=256i(e^2+2)+2048\).
All displayed divisions are fixed integer coefficients.

For all these choices, \(V_1>\lambda\sqrt{a_1}\) holds unconditionally.
Here is a finite proof rather than an asymptotic series argument.  Put
\[
 t_r=\frac{1024D^6-512D^5-128D^4-
                 \sum_{j=1}^r c_jD^{4-j}}{s_r}.
\]
The exact polynomial
\(t_r^2-(1024/s_r)^2D^{11}(D-1)\) has only positive nonzero coefficients,
and \(t_r>0\).  The remainders in \(N_*=D^jN_j+r_j\) are
\begin{align*}
 r_1&=3-A,&r_2&=3CD+3-A,\\
 r_3&=2e^2(A-1)D^2+r_2,&r_4&=2e^2(1-A-C)D^3+r_3.
\end{align*}
For \(A\ge3,C,i\ge1\), \(eA/D<9i/2\), \(eC/D<3i/2\), and
\(r_2,r_3>0\).  Thus orders one through three have adverse remainder
less than \(288i/s_r\).  For order four it is less than
\(252ie^2+(765/2)i\).  Each correction exceeds this loss plus \(\lambda\).
Finally, with \(X=e^2D^3(D-1)\), strict square-root concavity yields
\[
 \lambda\sqrt{DF}G<\frac\lambda2 eD^2\sqrt{1-1/D}\,N_*+\lambda v,
 \qquad v=\frac{CD+1}{2eD\sqrt{1-1/D}}<1.
\]
This proves \(T>\lambda\sqrt{DF}G\), and hence the required majorant.
No parity condition on \(g,h\) is needed.

\subsection{The equation and its equivalence}
The delivered single polynomial is
\[
 \boxed{\quad
 [\Theta^2-\Omega^2(a_0+\lambda^2a_1)]^2
           -4\lambda^2\Omega^4a_0a_1=0,\qquad
 \Theta=E_0-\Omega(V_0+V_1).
 \quad}
\]
All intermediate names are abbreviations, not extra unknowns.
Unconditionally
\(E_0=Jn+(J-2)W_1^2+3W_1C+2>0\).
A zero cannot have \(m=0\).  Real factorization gives signed roots
\(u_0^2=a_0,u_1^2=\lambda^2a_1\) and
\[
 n+W_1^2+Q/J=m(V_0+V_1+u_0+u_1).
\]
Both strict majorants make the last parenthesis positive, so \(m>0\)
follows before any assertion about rational roots.  The recovered carrier
range implies \(a_1>C^2A^4>K^2P^2>a_0\).  Consequently
\(x=u_0+u_1=\Theta/\Omega\ne0\); the identity
\(2xu_0=x^2+a_0-\lambda^2a_1\) proves both roots rational, hence integral.
Also \(u_1/\lambda\) is integral.  The factor identity now forces
\(J\mid Q\), proving soundness by the inherited ratio theorem.

Conversely use the ratio witnesses, choose positive roots, and set
\(n=m(V_0+V_1+u_0+u_1)-W_1^2-Q/J\).
It is an integer greater than zero because
\(W_1^2+Q/J<2D<V_1\).  This proves completeness and the claimed
consistency equivalence.  Natural-domain circuits use the bijective shifts,
not an unjustified alteration of the CRT index.

\subsection{Exact cancellation checks}
Writing \(\delta_0=V_0^2-a_0\),
\(\delta_1=V_1^2-\lambda^2a_1\), and
\(B_*=E_0^2-2E_0\Omega(V_0+V_1)+\Omega^2(\delta_0+\delta_1)\),
the two-defect identity is
\[
 F_r=B_*^2+4\Omega^2V_0V_1B_*+
             4\Omega^4(\delta_0V_1^2+a_0\delta_1).
\]
Finite remainder identities resolve \(D^8\delta_1\), giving a
division-free certificate equal to \(D^{16}F_r\).
The exact degree follows from a multivariate upper bound and a positive
nonzero diagonal leading monomial in the actual fixed coefficients.
Independent modular leading-coefficient checks use all four materialized
payloads.  The primary equation never divides by \(D\); its degree is
recovered by subtracting \(16\deg D\) from the certificate degree.
The package contains the explicit circuits, full proof, payload hashes,
symbolic ring checks, source checks, and independent literal replay.
These are mathematical proofs and reproducible finite audits, not a new
end-to-end kernel certificate or a minimality claim.
